\section{Industrial context and intended use}
\label{sec:context}

The system was motivated by paper-based production records in a bamboo-processing setting, where form values can be adjacent to employee identity, output accounting, and payroll decisions. The intended user is a reviewer who receives a digitized form, inspects machine candidates alongside image evidence, confirms or corrects values, and exports only confirmed records. The prototype demonstrates that this workflow can be represented in software, but the paper does not claim a measured deployment effect on cost, time, accuracy, wages, or organizational outcomes.

No production-form corpus was available to this study under an authorized research-use protocol. We therefore used deterministic synthetic forms for controlled and reproducible mechanism evaluation. Privacy, commercially sensitive content, and authorization requirements further constrain the acquisition and public release of such data. Synthetic forms preserve known labels, controlled perturbations, and distributable evidence. They do not reproduce handwriting variation, capture habits, paper wear, device diversity, or reviewer behavior in the plant.

For production adoption, the authority model should be embedded in a broader control system: least-privilege host access, encrypted storage and backups, retention schedules, separation of reviewer and payroll roles where required, audit review, incident response, and periodic sampling of confirmed records. Candidate--fact separation reduces one route by which machine error becomes operational truth; it does not replace these organizational controls.
